% =====================================================================
%  3. Variables cibles et stratégie de prédiction
% =====================================================================
\clearpage
\section{Variables cibles et stratégie de prédiction}
\label{sec:cibles}

\subsection{Quelles variables représentent la qualité ?}

À partir de l'EDA, nous retenons comme \textbf{cibles principales} les trois
propriétés de traction --- \YS, \UTS{} et allongement --- et comme \textbf{cible
exploratoire} l'énergie Charpy à température d'essai connue. Le
\cref{tab:cibles} justifie ce choix et donne les effectifs à chaque étape du
filtrage décrit en \cref{sec:prep}.

\begin{table}[htbp]
\centering
\small
\caption{Cibles retenues et effectifs de labels. \og Cohorte \fg{} désigne la
population de modélisation de 1\,555 lignes (\cref{subsec:cohorte}) ;
\og Groupes \fg{} le nombre de séries inférées portant au moins un label.}
\label{tab:cibles}
\begin{tabularx}{\textwidth}{@{}l l r r r r L@{}}
\toprule
Cible & Unité & Brut & Dédupl. & Cohorte & Groupes & Rôle et justification \\
\midrule
\YS           & MPa & 780 & 777 & 706 & 40 & Principale : propriété de dimensionnement, bonne couverture \\
\UTS          & MPa & 738 & 735 & 677 & 39 & Principale : résistance, complémentaire de \YS \\
Allongement   & \%  & 700 & 697 & 631 & 34 & Principale : ductilité, opposée à la résistance \\
Énergie \charpy & J & 879 & 877 & 861 & 24 & Exploratoire : protocoles hétérogènes (\cref{subsec:charpy_strategie}) \\
\midrule
Réduction de section & \% & 705 & 702 & 635 & -- & Complémentaire, non entraînée \\
Dureté        & kg/mm$^2$ & 138 & -- & -- & -- & Écartée : 80 valeurs exactes, protocoles Vickers mélangés \\
FATT50        & ?   & 31  & -- & -- & -- & Écartée : 31 valeurs, unité non documentée \\
\bottomrule
\end{tabularx}
\end{table}

Les distributions des cibles dans la cohorte sont résumées au \cref{tab:stats_cibles}.

\begin{table}[htbp]
\centering
\small
\caption{Statistiques des cibles dans la cohorte de modélisation.}
\label{tab:stats_cibles}
\begin{tabular}{@{}lrrrr@{}}
\toprule
Cible & Labels & Médiane & Minimum & Maximum \\
\midrule
\YS{} (MPa)              & 706 & 496,0 & 315,0 & 920,0 \\
\UTS{} (MPa)             & 677 & 576,0 & 447,0 & 1\,151,0 \\
Allongement (\%)         & 631 & 26,6  & 10,6  & 36,4 \\
Énergie \charpy{} (J)    & 861 & 100,0 & 3,0   & 270,0 \\
\bottomrule
\end{tabular}
\end{table}

\subsection{Pourquoi des régressions séparées ?}

Dans la cohorte, seules \textbf{609 lignes} renseignent simultanément les trois
propriétés de traction et \textbf{141} les quatre cibles. Un modèle
multi-sorties exigeant les trois propriétés de traction perdrait 14\,\% des
labels de \YS ; exiger les quatre cibles ne conserverait que 16 à 22\,\% des
labels selon la cible. Il faudrait sinon imputer des cibles, ce qui reviendrait
à entraîner sur des valeurs inventées.

\begin{decision}
Chaque cible dispose de \textbf{son propre modèle et de ses propres lignes
d'entraînement} : on apprend la cible $t$ sur toutes les lignes qui renseignent
$t$, sans exiger les autres. Une cible absente reste inconnue et
\textbf{n'est jamais imputée}. Aucune propriété mécanique n'est utilisée comme
entrée pour en prédire une autre, car elle ne serait pas disponible au moment
d'estimer une soudure non encore testée.
\end{decision}

\subsection{Le cas particulier de l'énergie Charpy}
\label{subsec:charpy_strategie}

Deux formulations sont possibles pour la ténacité : prédire l'énergie $E$ à une
température $T$ donnée, ou prédire la température $T_{100\,\mathrm{J}}$ (ou
$T_{28\,\mathrm{J}}$) à laquelle l'énergie atteint un seuil. Nous retenons la
première, \textbf{$E(T)$ avec la température d'essai en entrée}, formulation
utilisée par des travaux de référence sur cette famille de données~\cite{bhadeshia_1995,cool_1996}.
La température n'est jamais reconstruite : seules les 861 lignes de la cohorte
où elle est connue sont utilisées.

\begin{precaution}
La documentation définit bien deux colonnes distinctes (température en °C,
énergie en J), mais pas le protocole de chaque ligne. Une partie des couples
correspond probablement à des températures de transition à 28 ou 100\,J
(\cref{fig:charpy}), une autre à des courbes $E(T)$ complètes. Aucun seuil
28/100\,J ne sert à exclure automatiquement des lignes : cette tâche est traitée
comme \textbf{exploratoire} et les sources ont été auditées a posteriori
(\cref{subsec:charpy_audit}).
\end{precaution}

\subsection{Entrées autorisées et entrées exclues}
\label{subsec:entrees_exclues}

Une variable n'est utilisée en entrée que si elle serait connue \emph{avant}
l'essai mécanique, dans un usage réaliste du modèle. Le \cref{tab:entrees_exclues}
résume cette règle.

\begin{table}[htbp]
\centering
\small
\caption{Entrées autorisées et exclues des modèles.}
\label{tab:entrees_exclues}
\begin{tabularx}{\textwidth}{@{}l L@{}}
\toprule
Statut & Variables et raison \\
\midrule
\textbf{Autorisées} & Composition (C, Si, Mn, S, P, Ni, Cr, Mo, O, N), courant, tension, apport thermique, température entre passes, température et durée PWHT, procédé ; pour Charpy uniquement : température d'essai. \\
\addlinespace
Exclues (fuite) & Autres propriétés mécaniques, dureté, FATT50 : mesurées après l'essai, indisponibles pour une soudure nouvelle. \\
Exclues (fuite) & Identifiant, série, groupe, numéro de pli, drapeaux d'éligibilité : ils renseigneraient sur la source et non sur la soudure. \\
Exclues (disponibilité) & Microstructure : observée après soudage, 89 à 98 valeurs par colonne seulement. \\
Exclues (fiabilité) & V, Cu, Co, W, Ti, Al, B, Nb, Sn, As, Sb : couverture faible, censures nombreuses ou unités douteuses. Cette exclusion ne signifie pas que ces éléments n'ont pas d'effet métallurgique. \\
Non utilisées & Type de courant et polarité : conservés dans la table maître (modalité \texttt{missing} si non rapportés), non retenus dans les schémas d'entrées. \\
\bottomrule
\end{tabularx}
\end{table}

\subsection{Stratégie d'ensemble}

La stratégie de prédiction se déroule en cinq temps :
\begin{enumerate}
  \item \textbf{Construire une population fiable} : nettoyage déterministe, déduplication, groupes de séries et cohorte avec peu d'entrées manquantes (\cref{sec:prep}).
  \item \textbf{Fixer une référence minimale} : prédire la médiane de $y$ du train, sans utiliser $\bm{x}$. Tout modèle doit faire mieux.
  \item \textbf{Comparer des familles de modèles vues en cours}, du linéaire régularisé (Ridge) aux méthodes non linéaires (SVR, forêt aléatoire, $k$NN), avec réglage des hyperparamètres par validation croisée \textbf{imbriquée} (\cref{sec:methodes,sec:protocole}).
  \item \textbf{Évaluer selon deux usages} : interpolation dans des séries connues (découpage aléatoire) et transfert à de nouvelles séries (découpage groupé).
  \item \textbf{Exploiter les lignes non étiquetées} avec une méthode semi-supervisée (COREG), comparée à son exact pendant supervisé.
\end{enumerate}

Les lignes de la cohorte sans la cible forment le \textbf{pool non étiqueté} de
cette cible : 849 lignes pour \YS, 878 pour \UTS{} et 924 pour l'allongement.
Pour Charpy, \textbf{aucune} ligne de la cohorte n'a une température d'essai
connue sans énergie : le pool naturel est vide, et aucune température n'est
fabriquée pour en créer un.

\begin{aretenir}
Quatre régressions indépendantes : \YS, \UTS, allongement (principales) et
énergie Charpy $E(T)$ (exploratoire). Les entrées sont limitées à ce qui est
connu avant l'essai ; aucune cible n'est imputée ni utilisée comme entrée. Les
lignes sans cible alimentent l'étude semi-supervisée pour les trois cibles de
traction.
\end{aretenir}
